% ChargeShield-FL --- DSN 2027 --- Threat Model
% Portata in LaTeX 2026-09-16 (Sprint 10zz+117) dallo skeleton .docx.
% I cinque scenari sono ripresi dal .docx (struttura e contenuto invariati:
% erano gia' corretti e onesti, incluse le non-coperture dichiarate).
%
% AGGIUNTA rispetto al .docx: la sottosezione Where the Adversary Observes
% (sec:threat-observation) piu' la figura fig:adversary-points. Motivo, su
% richiesta esplicita dell'utente: i tre posizionamenti DP erano spiegati
% solo come meccanica nei Preliminaries (DOVE lungo la pipeline), mai come
% CHI ha quel privilegio in un deployment reale. Senza quel passaggio i tre
% bracci sembrano uno sweep di configurazione invece che un argomento di
% sicurezza.
%
% PUNTO DI ONESTA' introdotto qui e non presente nel .docx: dp-fedavg come
% punto di osservazione e' STRETTAMENTE PIU' FORTE dello Scenario 1. Sotto
% dp-fedavg e local l'update grezzo non lascia mai il client non
% privatizzato, quindi nessun aggregatore honest-but-curious lo vede mai.
% Dichiararlo rafforza l'argomento (un nullo sotto un avversario piu'
% privilegiato del nostro threat model limita anche gli altri due) ed evita
% un'obiezione facile.

\section{Threat Model}
\label{sec:threat-model}

We formalise five scenarios along a Purdue-model-inspired hierarchy of
attacker capability. The results in this paper concern Scenario~1
exclusively. Every other scenario is either handled by a different consumer
of the same ML Plane substrate and validated independently (Scenario~2), or
genuinely out of scope (Scenarios~3--4), reported as such rather than
silently assumed away.

\subsection{Scenario 1 --- Honest-but-Curious Aggregator (evaluated)}
\label{sec:threat-s1}

The server/aggregator follows the FL protocol faithfully --- it never
deviates and never sends malformed global weights --- but may passively
record and analyse every update it legitimately receives. Clients are honest
by construction. This is the only scenario our Yeom, Shadow and LiRA results
characterise.

\subsection{Where the Adversary Observes}
\label{sec:threat-observation}

Scenario~1 fixes the adversary's \emph{behaviour} but not its \emph{vantage
point}, and in a federated deployment those are separate questions. The
update a client produces passes through a privatisation pipeline
(Equation~\ref{eq:pipeline}) before it is aggregated, and how much of that
pipeline has already been applied when the adversary reads the update
determines what is left to extract. Figure~\ref{fig:adversary-points} makes
the three positions explicit.

\textbf{A3 --- noised update ($\tilde{g}$).} The adversary sees the update
after clipping and after noise. This is what an honest-but-curious
aggregator receives under a local-DP deployment, and equally what a network
observer would see on the wire. It is the least privileged position and the
one that matches the deployment most operators have in mind when they
say that they use DP.

\textbf{A2 --- clipped update ($\bar{g}$).} The adversary sees the update
after clipping but \emph{before} noise. This is not a hypothetical: under a
central-DP deployment the aggregator adds the noise itself, so it
legitimately receives the pre-noise value. A2 is therefore still Scenario~1
--- an aggregator that follows the protocol exactly and merely looks at what
the protocol hands it.

\textbf{A1 --- raw update ($g$).} The adversary sees the update before any
protective transformation. Under our deployments the raw update never leaves
the client unprivatised, so no honest-but-curious aggregator ever holds it;
A1 models an adversary \textbf{strictly stronger than Scenario~1}, such as
an insider or a compromised process on the charging-station controller
itself, operating inside the client boundary.

We state that asymmetry rather than smoothing it over, because it is what
makes the design informative. A1 upper-bounds what \emph{any} adversary can
extract from this channel: it observes the quantity from which A2 and A3 are
derived by information-destroying transformations. Consequently, if A1
yields no membership signal, a null at A2 or A3 cannot be credited to the
noise having suppressed a signal that was otherwise present --- there was
none to suppress. Reporting the three as competing mechanisms, rather than
as ordered privileges, would discard exactly this argument.

\subsection{Scenario 2 (a different consumer) and Scenarios 3--4 (out of scope)}
\label{sec:threat-s2}

Scenario~2, a malicious FL client, is not simply out of scope. The ML Plane
substrate has two consumers with orthogonal threat models over the same
observability layer: the Privacy Auditor (confidentiality, Scenario~1, the
subject of this paper's results) and ByzantineDetector (integrity,
Scenario~2). Scenario~2 belongs to that other consumer, is validated
separately from the privacy claims here, and is not conflated with
Scenario~1 in our results. As of this submission, ByzantineDetector's
Krum, cosine-similarity and CUSUM mechanisms are implemented and unit-tested
but have not been exercised end-to-end against an active Byzantine client on
real FL training data --- a gap we disclose rather than imply coverage we
have not measured.

Scenario~3 (compromised network gateway) and Scenario~4 (external passive
eavesdropper, defeated by transport-layer security) remain genuinely out of
scope: neither is an ML Plane consumer, and both are conventional
network-security concerns orthogonal to Scenarios~1 and~2 alike.

\subsection{Scenario 5 --- Adaptive Attacker (discussion only)}
\label{sec:threat-s5}

We discuss, without evaluating, two open questions: an attack tailored to
this exact autoencoder/FedProx/ACN-Data setup rather than the
literature-standard Yeom/Shadow/LiRA (S5a), and evasion of the Privacy
Auditor's own scoring (S5b). We explicitly reject a resolution we considered
during development --- that ByzantineDetector's stability controls would
catch an adaptive privacy attacker --- because Scenario~1 defines clients as
honest by construction: any attacker that deviates from protocol to evade
the Auditor is, by definition, Scenario~2. We therefore state a dual-defence
architecture with two orthogonal domains, integrity via ByzantineDetector
(Scenario~2) and privacy via the Auditor and DP (Scenario~1), rather than
claim that either component covers both.
